\section{Harness disclosure}
\label{sec:appendix-harness}

% SOURCE OF TRUTH: ../docs/HARNESS_DISCLOSURE.md, which carries the record for
% all three backends this project has. Reproduced here for the one backend
% this study used. Where the two disagree, the docs file is correct.

A result reported without its harness is not
reproducible~\citep{harnessdisclosure2026}. This appendix discloses the
harness under the seven layers that paper requires: execution, tools,
context, scheduling, observability, verification, governance. The project
has three backends; this study used exactly one, and the other two are
disclosed separately in the repository, because folding them into one table
would be the error the disclosure exists to prevent.

\paragraph{Execution.} An OpenAI-compatible HTTP server, reached through
Ollama's native \pth{/api/chat}. The target is \texttt{\NUM{\studyModel}}, at
a context window of 16{,}384, output capped at 4{,}096 tokens,
\texttt{temperature=0}, one process at a time. \NUM{\studyReadings} arm calls
plus \NUM{\aaPairs} for the A/A control. The server is local, so the
notional cost is zero and the guard is call count and wall clock.

\paragraph{Tools.} None. No tool use, no retrieval, no MCP, no filesystem
access. Nothing reads a working directory, so the auto-memory and
planted-file concerns that apply to the CLI backend do not arise here. The
skill body is the only intervention, and ``the agent looked it up'' cannot
explain a difference between arms.

\paragraph{Context.} The system prompt is the whole of it: task framing, the
response-format contract, and the arm's body. There is no host prompt
underneath and no session history. An arm that appends a skill to a
pre-existing host prompt is therefore refused on this backend rather than
silently approximated, because with no host prompt to append to, that call
would be the isolated arm wearing another arm's label.

\paragraph{Scheduling.} Arms ran interleaved across two passes in chunks of
eight items (\texttt{chunk=8}, \texttt{passes=2}, item-major ordering), ensuring
that any server drift or temporal shifts were distributed evenly across arms.
The A/A pass ran on the shared placebo after the main study passes. Model
residency is pinned at \texttt{keep\_alive: 60m}, sent on every request by the
provider that makes them, for the reason \Cref{sec:method} gives. The pin,
the temperature, the concurrency limit, the request timeout and the endpoint
are properties of the provider source and are written into the run manifest
(\pth{run.json}).

\paragraph{Observability.} Every call writes a record. Its fields are the
item id, the template id, the seed, the arm kind, the SHA-256 of the body
the arm delivered (null for \texttt{off}), the answer key, the parsed
answer, the parse status, the response text, the token counts, the duration,
and the schema version. Two ambiguities in the arm-kind field are resolved
by other fields: the evolved arms both record \texttt{candidate} and are
told apart by the hash, and the A/A pass records \texttt{placebo} under the
control's own hash and is told apart by its file, which is where
\texttt{de figures} reads the arm from. The records are published beside the
run. The accuracies, $p$-values and adjusted $q$-values in this paper are
the ones the run's own \texttt{analysis.json} computed and published;
everything else is recomputed from the records by \texttt{de figures} at
build time. The exceptions, numbers typed rather than generated, are
enumerated in \texttt{main.tex}: figures from cited work, configuration
constants, the probe and search figures of \Cref{sec:ceiling} and
\Cref{sec:whatengines}, and the per-template skew beside
\Cref{tab:templates}'s informedness.

\paragraph{Verification.} The resolved model id is recorded per call from
the response's own \texttt{model} field, not from what was requested.
Ollama's \pth{/api/show} card can be read as an isolation receipt, and
\texttt{assert\_isolated} refuses a card carrying a \texttt{SYSTEM} prompt,
the local analogue of a planted instruction file. The study path does not
call it: the only callers of \texttt{isolation\_receipt} in the repository
are tests, so this run holds no isolation receipt. Unlike the initial pilot
where reasoning output was omitted, the seven-arm study records the full
reasoning chain in the \texttt{reasoning} field on every record, satisfying
the transcript-level disclosure standard of \citet{biderman2024lessons}.

\paragraph{Governance.} The run is checkpointed and resumable, keyed so that
a resume cannot silently re-run or silently skip. The A/A pass has its own
checkpoint for that reason. The published directory declares its answer key,
its corpus fingerprints, its arm hashes and the notebook entry carrying its
prediction, and a gate refuses the run if that entry's first commit is not a
git ancestor of the run's commit.

\paragraph{One thing this backend cannot give.} The hosted screen of
\Cref{sec:ceiling} ran through the same module against a free API tier,
which publishes no model card, so no isolation receipt is obtainable there.
The harness has a function that would return \emph{no receipt obtainable}
rather than \emph{isolation verified}; the screen path does not call it
either, and no record of the study carries a receipt field.
